
\subsubsection{2.1 Static Analysis Feedback}

Static analysis tools like pylint and mypy analyze code structure without execution, detecting type mismatches, undefined variables, unreachable code, and style violations. The Static Analysis as Feedback Loop approach demonstrates that providing pylint/mypy errors as feedback during iterative refinement reduces security issues from 40\% to 13\% over 10 iterations. However, this work does not compare against execution-based feedback.

\subsubsection{2.2 Execution Feedback}

Execution-based feedback uses test results---pass/fail outcomes, error messages, and output comparisons---to guide refinement. Self-Refine demonstrates that LLMs can iteratively improve outputs using feedback, achieving approximately 20\% improvement across diverse tasks including code generation. CodeRL and related approaches train models with execution signals as rewards.

\subsubsection{2.3 Combined Approaches}

Some recent work combines both feedback types. The Helping LLMs Improve Code Generation approach provides both testing results and static analysis in a unified feedback loop. However, this work does not decompose individual contributions or measure overlap between signal types.

\subsubsection{2.4 Research Gap}

Prior work establishes that both static and execution feedback improve code quality independently. What remains unknown is whether they detect the same or different errors. This study addresses this gap by directly measuring error class overlap.

